% === Chapter 7: Efficiency & Data Engineering ===
\section{Thread 7: Efficiency \& Data Engineering}
\label{sec:t7}

\begin{problembox}
The computational overhead of post-training (SFT, RLHF, DPO, etc.) is prohibitively high:
full fine-tuning a 175B-parameter model requires TB-scale GPU memory and hundreds of GPUs;
the PPO phase of RLHF demands maintaining four large model replicas simultaneously---the policy model, reference model, reward model, and value model;
high-quality human-annotated data is scarce and expensive to obtain.
These cost barriers severely constrain both the research iteration speed and the deployment coverage of post-training techniques.
\textbf{How can we achieve comparable or superior post-training results with lower computational cost and less data?}
\end{problembox}

\noindent
This question has given rise to five major research directions in Thread~7:
(1) parameter-efficient fine-tuning, which approximates full fine-tuning performance with a minimal number of trainable parameters;
(2) data selection \& curation, which identifies the most valuable subsets from massive instruction datasets;
(3) synthetic data generation, which leverages strong models to produce training data at scale;
(4) distillation for alignment, which transfers alignment capabilities from large models to smaller ones;
(5) compute-efficient RL, which simplifies the training architecture of RLHF.
These directions are closely synergistic---for example, QLoRA makes it possible to perform RLHF on consumer-grade GPUs, while synthetic data provides low-cost training corpora for LoRA fine-tuning.

\subsection{Problem Origin: The Cost Bottleneck of Post-Training}
\label{sec:t7-origin}

The cost bottleneck of post-training can be understood along three dimensions.

\paragraph{Computational cost.}
Full fine-tuning requires maintaining gradients and optimizer states for all model parameters.
Taking the Adam optimizer as an example, each parameter requires storing three additional tensors: the gradient (FP16), the first moment (FP32), and the second moment (FP32).
This means that the optimizer states alone for a 175B-parameter model require approximately 1.2TB of GPU memory.
RLHF further amplifies this problem:
the PPO training phase requires simultaneously loading the policy model, the reference model (a frozen copy), the reward model, and the value model---the memory footprint of four large models makes even a 7B-scale model require multiple A100 GPUs \cite{santacroce2023efficient}.
\crossref{2}{sec:t2}

\paragraph{Data cost.}
The human annotation cost for high-quality instruction-following data is substantial:
InstructGPT \cite{ouyang2022instructgpt} employed 40 annotators, and the labeling process involved complex quality control procedures.
Preference data annotation is even more challenging---annotators must compare the quality of two model outputs, a task for which inter-annotator agreement tends to be low.
\crossref{1}{sec:t1}

\paragraph{Iteration speed.}
The compounding of the above computational and data costs results in extremely long experiment cycles.
A single RLHF training experiment may take days or even weeks, severely constraining the efficiency of hyperparameter search and methodology iteration.
This constitutes a nearly insurmountable barrier, particularly for academic researchers and small teams.

\evolution{Full fine-tuning 175B: requires 1.2TB memory, hundreds of GPUs}{LoRA: trainable parameters reduced to 0.01\%, single-GPU training}{Parameter efficiency}

\subsection{Foundational Methods: LoRA and QLoRA}
\label{sec:t7-foundation}

\subsubsection{LoRA: Low-Rank Adaptation}
\label{sec:t7-lora}

\landmark{hu2022lora} introduced a key insight:
\textbf{the weight updates during adaptation of pre-trained language models exhibit low-rank structure}.
This implies that although the number of model parameters is enormous, the ``information content'' required for fine-tuning is far smaller than the full parameter space of the model.

\paragraph{Core method.}
For a pre-trained weight matrix $W_0 \in \mathbb{R}^{d \times k}$,
LoRA constrains its update to a low-rank decomposition:
\begin{equation}
h = W_0 x + \Delta W x = W_0 x + B A x
\label{eq:lora}
\end{equation}
where $B \in \mathbb{R}^{d \times r}$, $A \in \mathbb{R}^{r \times k}$,
and rank $r \ll \min(d, k)$.
During training, the original weights $W_0$ are completely frozen (receiving no gradient updates),
and only $A$ and $B$ are trainable parameters.
$A$ is initialized with random Gaussian values, while $B$ is initialized to the zero matrix,
ensuring that $\Delta W = BA = 0$ at the start of training---i.e., the adapted model begins in a state identical to the pre-trained model.
In practice, a scaling factor $\alpha / r$ is also introduced so that switching ranks does not require significant learning rate adjustments.

\paragraph{Key experimental findings.}
Experiments on GPT-3 175B demonstrated the following:
\begin{itemize}[leftmargin=2em]
  \item \textbf{Parameter efficiency}:
    LoRA requires only 4.7M trainable parameters (approximately 0.003\% of the 175B total),
    reducing the checkpoint size from 350GB to 35MB and memory requirements from approximately 1.2TB to approximately 350GB.
  \item \textbf{Performance}:
    On tasks such as WikiSQL (+0.4\%) and SAMSum (+0.7\%),
    LoRA matched or even surpassed full fine-tuning performance,
    and outperformed prior parameter-efficient methods:
    the \textbf{Adapter} approach proposed by Houlsby et al.\ \cite{houlsby2019adapters}, which inserts small bottleneck layers into each Transformer layer,
    and \textbf{Prefix Tuning} by Li \& Liang \cite{li2021prefix}, which freezes model parameters and only optimizes continuous prefix embeddings at the input.
    LoRA outperformed both in terms of performance, and more crucially, introduced no additional inference latency.
  \item \textbf{Zero inference latency}:
    At deployment time, $\Delta W = BA$ can be merged back into $W_0$,
    yielding a weight matrix $W' = W_0 + BA$ with the same structure as the original model,
    introducing no additional inference computation or latency.
    This property is a significant advantage of LoRA over adapter and prefix tuning approaches.
  \item \textbf{Rank selection}:
    Experiments revealed that extremely low ranks ($r = 1$ or $r = 2$) are sufficient for many tasks,
    indicating that the intrinsic dimensionality of adaptation is far lower than the model parameter space.
    Applying LoRA to $W_q$ and $W_v$ (the query and value projection matrices in attention)
    with $r = 4$ yielded the best results.
\end{itemize}

\paragraph{Why does low rank work?}
Through their analysis, Hu et al.\ found that
the subspace overlap between $\Delta W$ and $W_0$ is very low---that is, the directions required for adaptation primarily reside in subspaces that were not heavily emphasized during pre-training.
This explains why extremely low-rank updates can effectively alter model behavior:
the essence of post-training is not to rewrite all of the model's knowledge,
but rather to ``modulate'' how the model expresses its existing capabilities within a low-dimensional subspace.
This is highly consistent with the Superficial Alignment Hypothesis proposed by LIMA \cite{zhou2023lima}:
\textbf{the amount of information required for alignment is far less than the knowledge imparted to the model during pre-training}.
\crossref{1}{sec:t1}

\begin{solutionbox}
The core contribution of LoRA is transforming fine-tuning from a full parameter-space problem into a low-rank subspace problem:
freezing the pre-trained weights $W_0$ and learning only the low-rank update $\Delta W = BA$,
achieving 10,000$\times$ parameter compression and 3$\times$ memory savings,
with zero additional inference latency at deployment.
\end{solutionbox}

\subsubsection{QLoRA: Fusing Quantization with LoRA}
\label{sec:t7-qlora}

\landmark{dettmers2023qlora} combined quantization techniques with LoRA,
further compressing memory requirements to a fraction of what LoRA alone needs.
The core objective was: \textbf{to fine-tune a 65B-parameter model on a single 48GB GPU},
with no performance degradation relative to 16-bit full fine-tuning.

\paragraph{Three key techniques.}
QLoRA introduced three complementary innovations:

\begin{enumerate}[leftmargin=2em]
  \item \textbf{4-bit NormalFloat (NF4)}:
    An information-theoretically optimal quantization data type.
    Its design is based on a key observation:
    pre-trained neural network weights approximately follow a zero-mean normal distribution $\mathcal{N}(0, \sigma^2)$.
    NF4 sets quantization anchor points by computing equal-probability quantiles of the standard normal distribution:
    \begin{equation}
    q_i = \frac{1}{2}\left(Q_X\!\left(\frac{i}{2^k+1}\right) + Q_X\!\left(\frac{i+1}{2^k+1}\right)\right)
    \end{equation}
    where $Q_X$ is the quantile function of the standard normal distribution and $k = 4$ is the number of quantization bits.
    Each anchor interval contains an equal probability mass of weight values,
    thereby maximizing the information carried by each quantization bit in an information-theoretic sense.
    Experiments showed that NF4 consistently outperforms FP4 and Int4 in quantization error.

  \item \textbf{Double Quantization (DQ)}:
    The quantization constants generated during quantization themselves consume memory---every 64 parameters correspond to one FP32 scaling factor, amounting to an additional 0.5 bits per parameter.
    DQ applies a second round of quantization to these quantization constants (using 8-bit quantization with a larger block size of 256),
    reducing the overhead from 0.5 bits per parameter to approximately 0.127 bits,
    saving approximately 3GB of memory on a 65B model.

  \item \textbf{Paged Optimizers}:
    Leveraging NVIDIA's unified memory mechanism,
    optimizer states are automatically paged to CPU memory when GPU memory is insufficient,
    preventing out-of-memory errors caused by uneven sequence lengths within mini-batches.
    This is particularly important for handling memory spikes during the generation phase of RLHF.
\end{enumerate}

\paragraph{Complete computation flow.}
The forward pass of QLoRA can be expressed as:
\begin{equation}
Y^{\text{BF16}} = X^{\text{BF16}} \cdot \text{doubleDequant}(c_1^{\text{FP32}}, c_2^{\text{k-bit}}, W^{\text{NF4}}) + X^{\text{BF16}} \cdot L_1^{\text{BF16}} \cdot L_2^{\text{BF16}}
\end{equation}
where the base model weights are stored in NF4 (dequantized to BF16 during the forward pass),
while the LoRA adapters $L_1$ and $L_2$ are always computed in BF16.
Gradients are backpropagated only through the LoRA adapters; the base weights remain frozen throughout training.

\paragraph{Guanaco: Large-scale validation of QLoRA.}
Based on QLoRA, Dettmers et al.\ trained the Guanaco family of models,
performing single-GPU fine-tuning on the OASST1 dataset (approximately 9,000 multi-turn dialogues).
Key results:
\begin{itemize}[leftmargin=2em]
  \item Guanaco 65B achieved 99.3\% of ChatGPT's performance on the Vicuna benchmark,
    requiring only a single 48GB GPU and 24 hours of training.
  \item Guanaco 33B surpassed all prior open-source models with just 12 hours of training on a single 24GB GPU.
  \item \textbf{4-bit QLoRA matches the performance of 16-bit full fine-tuning}:
    controlled experiments across multiple tasks and model scales demonstrated that
    QLoRA introduces no systematic performance degradation.
\end{itemize}

\begin{solutionbox}
QLoRA achieves fine-tuning of a 65B model on a single 48GB GPU through
NF4 quantization (information-theoretically optimal) + Double Quantization (compressing quantization constants)
+ Paged Optimizers (handling memory spikes),
with experiments confirming that 4-bit QLoRA fully matches the performance of 16-bit full fine-tuning.
\end{solutionbox}

\subsection{Limitations and Branches: Beyond Parameter Efficiency}
\label{sec:t7-branches}

LoRA and QLoRA addressed the parameter efficiency problem, but the efficiency challenges of post-training extend far beyond this.
Parameter efficiency is only one dimension of computational cost---data acquisition and selection, training data quality and diversity, distillation efficiency, and the architectural complexity of RL training are all independent and significant efficiency bottlenecks.
Consequently, Thread~7 diverged into five parallel research branches following the foundational methods.

\subsubsection{Branch A: Subsequent Developments in Parameter-Efficient Fine-Tuning}
\label{sec:t7-peft}

After LoRA established the low-rank adaptation paradigm, subsequent work improved upon it from multiple angles.

\paragraph{DoRA: Weight-Decomposed Low-Rank Adaptation.}
Liu et al.\ \cite{liu2024dora} observed significant differences between full fine-tuning and LoRA in the magnitude and direction change patterns of weight updates.
DoRA decomposes the weight matrix into magnitude and direction components: $W = m \cdot \frac{V}{\|V\|}$,
where the magnitude $m$ is learned independently and the direction $V$ is adjusted through LoRA's low-rank updates.
This decomposition enables DoRA to consistently outperform LoRA across various NLP and vision tasks without introducing additional inference overhead (weights can likewise be merged at deployment);
see \crossref{5}{sec:t5} for post-training methods on vision-language models.

\paragraph{PERL: Parameter-Efficient RLHF.}
Sidahmed et al.\ \cite{sidahmed2024perl} systematically applied LoRA to the PPO phase of RLHF,
demonstrating that using LoRA for both the reward model and the policy model can reduce the memory requirements of RLHF by over 50\%, with final alignment performance comparable to full fine-tuning RLHF.
This result shows that the combination of parameter-efficient methods with RL-based alignment is viable.
\crossref{2}{sec:t2}

\paragraph{DoRAN: Stabilized DoRA.}
Diep et al.\ \cite{diep2025doran} proposed \textbf{DoRAN},
which introduces noise injection and auxiliary networks on top of DoRA's magnitude-direction decomposition,
addressing DoRA's instability issues under certain training configurations.
DoRAN significantly improves training robustness while preserving DoRA's accuracy advantages, performing particularly well in RL-based post-training scenarios.

\subsubsection{Branch B: Data Selection \& Curation}
\label{sec:t7-data-selection}

\begin{problembox}
The effectiveness of instruction tuning is highly dependent on data quality, yet high-quality data annotation is expensive.
Given large-scale automatically generated or web-crawled instruction datasets,
how can we select the most valuable subsets to maximize training efficiency?
\end{problembox}

\paragraph{LIMA: The Superficial Alignment Hypothesis.}
\landmark{zhou2023lima} fine-tuned LLaMA 65B using only 1,000 carefully curated instruction-response pairs,
and the resulting model was preferred over or tied with DaVinci003 in 43\% of cases in human evaluation.
The resulting \textbf{Superficial Alignment Hypothesis}---that alignment merely teaches the model to present its existing capabilities in the correct format---fundamentally reframed the optimization objective of data engineering:
the key factor is not data volume, but rather \textbf{data quality and diversity}.
Detailed experimental results and an in-depth discussion of the hypothesis can be found in Thread~1 (\S\ref{sec:t1-quality}).

\paragraph{Subsequent data selection methods.}
LIMA's findings inspired a series of automated data selection methods:

\begin{itemize}[leftmargin=2em]
  \item \textbf{AlpaGasus} \cite{chen2024alpagasus}:
    Used ChatGPT as a quality evaluator to score Alpaca's 52K data samples,
    retaining only 9K high-quality samples for training, which outperformed training on the full 52K dataset.

  \item \textbf{DEITA} \cite{liu2024deita}:
    Proposed a three-dimensional data evaluation framework---complexity (instruction complexity),
    quality (response quality), and diversity (sample diversity).
    Through scoring and diversity-aware sampling, selecting only 6K samples matched the performance of 10K randomly sampled examples.

  \item \textbf{IFD (Instruction Following Difficulty)} \cite{li2023ifd}:
    Used the model's own perplexity changes as a proxy for data value---if a sample's IFD score is high (i.e., the model's perplexity changes substantially before and after incorporating the instruction),
    it indicates that the sample contains information the model has not yet mastered, making it more valuable for training.

  \item \textbf{Superfiltering} \cite{li2024superfiltering}:
    Discovered that IFD rankings produced by small models (e.g., GPT-2) are highly consistent with those from large models,
    enabling the use of small models for preliminary filtering and substantially reducing the computational cost of the selection stage.

  \item \textbf{LESS} \cite{xia2024less}:
    Based on the theoretical framework of influence functions,
    selects the most relevant training samples according to the validation loss gradient on the target task.
    Only 5\% of the data is needed to match or even exceed the performance achieved with the full dataset.
\end{itemize}

\subsubsection{Branch C: Synthetic Data Generation}
\label{sec:t7-synthetic}

\begin{problembox}
Human-annotated instruction data is high in quality but expensive to obtain and difficult to scale.
Can existing strong models (e.g., GPT-4) be leveraged to automatically generate large-scale, high-quality training data?
\end{problembox}

\paragraph{Self-Instruct: Generating data from the model itself.}
\landmark{wang2023selfinstruct} proposed a bootstrapping pipeline:
starting from a small set (175) of manually written seed instructions,
the dataset is iteratively expanded through four steps:
(1) having the model generate new task instructions,
(2) determining whether the new instruction is a classification task,
(3) generating corresponding input instances,
(4) generating outputs.
Through ROUGE-L deduplication and heuristic filtering,
52K instruction data samples were ultimately generated from GPT-3 (text-davinci-001).
Fine-tuning GPT-3 on this self-generated data substantially outperformed the original GPT-3 in user-oriented evaluation,
approaching the level of InstructGPT.
Self-Instruct directly gave rise to open-source projects such as Stanford Alpaca \cite{taori2023alpaca},
establishing the ``training LLMs with LLMs'' data production paradigm.

\paragraph{Orca: Explanation Tuning.}
\landmark{mukherjee2023orca} identified a key limitation of Self-Instruct and Alpaca-style methods:
they only distill the teacher model's final outputs,
discarding the reasoning process behind those outputs.
Orca proposed \textbf{explanation tuning}---using carefully designed system instructions (such as ``explain step by step,'' ``think like you are answering to a five year old,'' and 16 other diverse reasoning prompts)
to guide GPT-4 to produce not only answers but also detailed reasoning processes (explanation traces).
Training data was sourced from 5M samples from FLAN-v2 (with ChatGPT-generated traces) plus 1M samples (with GPT-4-generated traces),
employing a \textbf{progressive learning} strategy:
training first on the easier ChatGPT traces, then continuing on the more challenging GPT-4 traces.
Orca 13B achieved ChatGPT-level performance on BigBench Hard,
demonstrating that \textbf{distilling reasoning processes} is more effective than merely distilling outputs.
The subsequent Orca 2 \cite{mitra2023orca2} further enabled small models to adaptively select reasoning strategies (e.g., step-by-step, direct answer) based on task difficulty,
surpassing larger-scale models on multiple reasoning benchmarks.
\crossref{4}{sec:t4}

\paragraph{Phi-1: Textbook-quality data.}
\landmark{gunasekar2023phi1} validated the decisive role of data quality from a unique angle.
The direct precursor to this work was TinyStories \cite{eldan2023tinystories}---Eldan \& Li demonstrated that extremely small models ($<$10M parameters) trained solely on short children's stories generated by GPT-3.5/GPT-4 could fluently generate coherent English, providing the first quantitative evidence of the decisive impact of synthetic data quality on small model capabilities.
Inspired by this, Gunasekar et al.\ did not filter from real web data
but instead used GPT-3.5 to directly generate ``textbook quality'' synthetic data---structured, progressive, and comprehensive textbook-grade code training data that covers core concepts.
A model with only 1.3B parameters achieved 50.6\% pass@1 on HumanEval,
surpassing many models more than ten times its size.
Phi-1.5 \cite{li2023phi15} extended this approach to the general NLP domain,
where a 1.3B model trained on synthetic textbook data approached the performance of 5$\times$--10$\times$ larger models on reasoning benchmarks,
further demonstrating that \textbf{data quality can compensate for model scale}.

\paragraph{WizardLM: Evol-Instruct.}
\landmark{xu2024wizardlm} proposed the Evol-Instruct method---performing multiple rounds of ``evolutionary'' rewriting on initial instructions,
progressively increasing instruction complexity and diversity by adding constraints, introducing multi-step reasoning, and concretizing scenarios.
This data augmentation strategy does not require re-annotation by humans
but instead leverages the LLM itself to systematically complexify existing instructions.
WizardLM 7B approached ChatGPT-level performance in human evaluation,
demonstrating a pathway to improving model capabilities by enhancing data \textbf{complexity} (rather than quantity).

\paragraph{UltraChat and UltraFeedback: Large-scale synthetic datasets.}
UltraChat \cite{ding2023ultrachat} constructed approximately 1.5M multi-turn dialogue data samples covering diverse topics including world knowledge, creative writing, and data analysis,
providing large-scale, diversified training resources for SFT.
UltraFeedback \cite{cui2023ultrafeedback} targeted preference learning---collecting responses from multiple models to the same prompt and using GPT-4 for preference annotation,
generating large-scale preference pair data.
Together, these two datasets form the data foundation of the Zephyr \cite{tunstall2023zephyr} pipeline.

\paragraph{Self-Alignment with Instruction Backtranslation (Humpback).}
Li et al.\ \cite{li2023humpback} proposed a reverse data generation method:
first sampling high-quality text from the web corpus as ``responses,''
then training a model to reverse-generate instructions for these responses (backtranslation),
followed by quality filtering of the generated (instruction, response) pairs.
After two iterations,
Humpback surpassed all open-source models trained with non-distillation methods on AlpacaEval.

\paragraph{RLAIF vs.\ RLHF: Substituting AI feedback for human feedback.}
\landmark{lee2023rlaif} systematically compared the effectiveness of AI-generated preference annotations (RLAIF) with human annotations (RLHF).
The core finding was that
on tasks such as summarization and helpful dialogue,
\textbf{the final model performance of RLAIF shows no significant difference from RLHF}---AI-annotated preference data can effectively substitute for human annotations.
This conclusion provides theoretical support for fully synthetic alignment pipelines
and lays the foundation for zero-human-annotation methods such as Zephyr.
\crossref{2}{sec:t2}

\subsubsection{Branch D: Distillation for Alignment}
\label{sec:t7-distillation}

\begin{problembox}
Large models (e.g., GPT-4) deliver exceptional performance but are prohibitively expensive to deploy.
Can the alignment capabilities of large models be ``distilled'' into smaller models,
enabling low-cost deployment of small models while preserving the alignment quality of large models?
\end{problembox}

\paragraph{Zephyr: Direct Distillation of LM Alignment.}
\landmark{tunstall2023zephyr} proposed a three-stage distillation pipeline that
\textbf{requires no human annotation whatsoever} to transfer alignment capabilities from large models to a 7B model:

\begin{enumerate}[leftmargin=2em]
  \item \textbf{Distilled SFT (dSFT)}:
    SFT was performed on the Mistral-7B \cite{jiang2023mistral} base model using the UltraChat dataset (approximately 200K multi-turn dialogues generated by ChatGPT).
    Mistral-7B itself is a pivotal base model---through architectural optimizations such as sliding window attention and grouped-query attention,
    it surpassed LLaMA-2 13B on multiple benchmarks with only 7B parameters,
    providing a high-quality training starting point for Zephyr.
    The data underwent deduplication and quality filtering to ensure diversity and quality in the training corpus.

  \item \textbf{AI Feedback (AIF)}:
    Using the UltraFeedback dataset---for each prompt, responses from 4 different models were collected,
    and GPT-4 scored the response quality (on a 1--10 scale).
    The highest-scored and lowest-scored responses were paired to form preference pairs.
    This step was entirely performed by AI, with \textbf{zero human annotation}.

  \item \textbf{Distilled DPO (dDPO)}:
    The dSFT model was trained with DPO using the above AI Feedback data.
    The dSFT checkpoint served as the reference model for DPO,
    training on approximately 64K preference pairs.
\end{enumerate}

\noindent
Zephyr-7B-$\beta$ achieved a score of 7.34 on MT-Bench,
\textbf{surpassing LLaMA2-Chat-70B} (6.86),
and approaching GPT-3.5-turbo (7.94)---this was the first time a 7B model outperformed a 70B RLHF model on a dialogue quality benchmark.

\paragraph{Key ablation experiments.}
The ablation experiments of Zephyr revealed several important findings:
\begin{itemize}[leftmargin=2em]
  \item There was a substantial gap between dSFT alone without dDPO (score 5.32) and the full pipeline (7.34),
    indicating that \textbf{preference optimization is indispensable},
    even when the data comes from distillation.
  \item Applying DPO directly to the base model without dSFT (score 4.76) performed poorly,
    indicating that the SFT stage provides an essential warm start.
  \item During DPO training, the training accuracy converged to nearly 1.0 (i.e., the model fully fit the preference data),
    yet the MT-Bench score \textbf{did not decline}---this ``overfitting without degradation'' phenomenon suggests
    that DPO exhibits good generalization under distribution-shifted evaluation.
\end{itemize}

\paragraph{Limitations of distillation: The False Promise.}
However, Gudibande et al.\ \cite{gudibande2023false} raised a pointed critique of the distillation paradigm.
They found that small models trained by imitating ChatGPT outputs (e.g., Alpaca-family models),
while quickly approaching the teacher in \textbf{style} (fluency, formatting),
showed no significant narrowing of the gap in \textbf{factuality} and \textbf{reasoning ability}---imitation primarily learned ``how to speak'' rather than ``how to think.''
This finding provides corroborating evidence for Orca's explanation tuning from the negative side:
output-level imitation alone is indeed insufficient;
process-level imitation (distilling reasoning processes) is necessary to transfer deeper capabilities.

\paragraph{Other distillation methods.}
\begin{itemize}[leftmargin=2em]
  \item \textbf{MiniLLM} \cite{gu2023minillm}:
    Replaced standard forward KL divergence with reverse KL divergence,
    training the student model to better fit the high-probability regions of the teacher distribution,
    consistently outperforming standard distillation on text generation tasks.

  \item \textbf{GKD (Generalized Knowledge Distillation)} \cite{agarwal2023gkd}:
    Proposed on-policy distillation---computing the KL divergence loss on the student model's own output distribution
    rather than on the teacher's output distribution.
    This alleviates the distribution shift between train-time (teacher distribution) and test-time (student distribution).

  \item \textbf{Lion} \cite{jiang2023lion}:
    Introduced an adversarial training framework,
    alternately training the student model (to imitate the teacher) and
    a discriminator (to distinguish student outputs from teacher outputs),
    improving distillation quality through adversarial signals.
\end{itemize}

\subsubsection{Branch E: Compute-Efficient Reinforcement Learning}
\label{sec:t7-compute-rl}

\begin{problembox}
PPO is the standard RL algorithm for RLHF, but its training architecture is extremely complex:
it requires simultaneously maintaining 4 models (policy, reference, reward, value),
incurs massive memory overhead, suffers from high training instability, and is sensitive to hyperparameters.
Can simpler RL algorithms be designed to replace PPO?
\end{problembox}

\paragraph{GRPO: Group Relative Policy Optimization.}
\landmark{shao2024grpo} proposed an elegant solution:
\textbf{completely removing the critic (value) model from PPO},
instead using the relative rewards among a group of outputs for the same prompt as the advantage baseline
(algorithmic details and the complete derivation are provided in \S\ref{sec:t2-grpo};
the special significance for reasoning scenarios is discussed in \S\ref{sec:t4-grpo}).

\paragraph{Efficiency gains.}
GRPO delivers threefold efficiency improvements:
\begin{itemize}[leftmargin=2em]
  \item \textbf{Elimination of the value model}:
    Memory is directly reduced by approximately 25\% (for equally-sized models),
    and the training architecture is simplified from 4 models to 3 (policy + reference + reward).
  \item \textbf{No per-token advantage estimation required}:
    PPO requires the value model to produce a value estimate for each token to compute GAE.
    GRPO directly uses outcome-level group-relative rewards,
    completely avoiding the error accumulation inherent in value estimation.
  \item \textbf{Natural compatibility with verifiable rewards}:
    For tasks with verifiable correctness such as mathematics and programming,
    rewards can be directly based on answer correctness (0/1 or rule-based),
    eliminating the need for a trained reward model.
    The final setup requires maintaining only 2 models (policy + reference),
    bringing the overhead close to that of DPO.
\end{itemize}

\noindent
GRPO achieved 51.7\% on MATH and 88.2\% on GSM8K with DeepSeekMath 7B,
significantly outperforming PPO and DPO baselines.
It was subsequently adopted by DeepSeek-R1 \cite{deepseek2025r1} as the core algorithm for large-scale reasoning training.
GRPO's success validates an efficiency principle:
\textbf{removing unnecessary components (the value model) not only reduces cost but can also improve performance}---because value estimation itself is an additional source of error.

\paragraph{Other compute-efficient RL methods.}
\begin{itemize}[leftmargin=2em]
  \item \textbf{Efficient RLHF} \cite{santacroce2023efficient}:
    Through engineering techniques such as gradient checkpointing,
    mixed-precision training, and distributed policy optimization,
    reduced PPO's memory usage by approximately 50\% without modifying the algorithm.

  \item \textbf{ReMax} \cite{li2023remax}:
    Replaced PPO with a variant of REINFORCE,
    using the sample mean as a baseline (similar in spirit to GRPO),
    eliminating the value model and GAE computation.
    Theoretical analysis of ReMax demonstrated that under the specific settings of LLM alignment,
    this simplification does not introduce excessive variance.

  \item \textbf{REINFORCE++} \cite{hu2025reinforcepp}:
    Built upon REINFORCE by introducing token-level KL penalties and PPO-style clipping,
    incorporating PPO's stability improvements while maintaining the simplicity of REINFORCE,
    achieving a favorable balance between performance and efficiency.

  \item \textbf{OpenRLHF} \cite{hu2024openrlhf}:
    An open-source RLHF training framework
    supporting multiple algorithms including PPO, GRPO, and REINFORCE++,
    with efficient distributed training orchestration via Ray.
\end{itemize}

\evolution{PPO: 4 models + GAE + complex training loop}{GRPO: 3 models + group-relative reward + simple objective}{Removal of the value model}

\subsection{Convergence and Frontiers: Scaling Laws for Post-Training}
\label{sec:t7-scaling}

After efficiency challenges were addressed at the technical level (parameter efficiency, data engineering, RL simplification),
a more fundamental question emerged:
\textbf{How does post-training performance scale with compute budget, data volume, and model size?}
Answering this question requires establishing scaling laws for the post-training domain.

\subsubsection{Data-Constrained Scaling Laws}
\label{sec:t7-scaling-data}

\landmark{muennighoff2023scaling} extended Chinchilla's scaling laws to
\textbf{data-constrained} scenarios.
Muennighoff et al.\ systematically studied the effects of data repetition (multi-epoch training)
across 400+ training experiments
(model scales from 10M to 9B, with training token counts reaching 900B).

\paragraph{Core findings.}
\begin{enumerate}[leftmargin=2em]
  \item \textbf{Effective data volume}:
    When unique data $U_D$ is repeated for $R_D$ epochs,
    the equivalent amount of ``new data'' can be described by the following formula:
    \begin{equation}
    D' = U_D + U_D \cdot R_D^* \cdot \left(1 - e^{-R_D / R_D^*}\right)
    \label{eq:effective-data}
    \end{equation}
    where $R_D^* \approx 15$ is the characteristic epoch count.
    Intuitively, the first 4 epochs of data repetition are approximately equivalent to using new data,
    but beyond approximately 16 epochs, the marginal benefit of each additional epoch diminishes sharply.

  \item \textbf{Optimal resource allocation under data constraints}:
    In data-constrained scenarios (i.e., when available unique data is limited),
    the optimal strategy is to allocate more compute to increasing the number of epochs
    rather than increasing model parameters.
    This stands in sharp contrast to the data-unconstrained scenario,
    where Chinchilla recommends proportional scaling of model and data.

  \item \textbf{Code mixing}:
    Incorporating code data into natural language training yields approximately 2$\times$
    effective token gains---the structured nature of code data provides useful inductive biases for natural language tasks.

  \item \textbf{Data deduplication vs.\ filtering}:
    On downstream task evaluation,
    perplexity-based filtering is more effective than deduplication---while deduplication intuitively reduces redundant information,
    its impact on final task performance is less pronounced than quality filtering.
\end{enumerate}

\paragraph{Implications for post-training.}
Although this work primarily studied pre-training,
its conclusions have direct relevance for post-training data engineering:
(1) in scenarios where instruction tuning data is limited, moderate multi-epoch training is justified;
(2) data augmentation (e.g., paraphrasing, backtranslation) can effectively increase data diversity,
breaking through the diminishing returns of epoch repetition;
(3) data quality filtering should take priority over simple deduplication.

\subsubsection{Post-Training-Specific Scaling Studies}
\label{sec:t7-scaling-pt}

\paragraph{Does RLHF Scale?}
Hou et al.\ \cite{hou2024doesrlhfscale} systematically investigated
the relationship between RLHF effectiveness and model scale as well as data volume.
Key findings include:
RLHF benefits increase with model scale (scaling with model size),
but the reward model overoptimization problem also intensifies---larger policy models become more adept at ``exploiting'' the reward model's weaknesses.
This implies an important scaling tension:
\textbf{larger models require stronger reward models for effective optimization guidance};
otherwise, reward hacking will offset the gains from scale.
\crossref{2}{sec:t2}

\paragraph{Scaling Laws for Reward Model Overoptimization.}
Rafailov et al.\ \cite{rafailov2024scaling} extended the scaling behavior of reward model overoptimization
to DPO and other direct alignment algorithms (DAAs).
The study found that DAAs also exhibit overoptimization---when KL divergence exceeds a certain threshold,
the true reward (gold standard evaluation) begins to decline
while the proxy reward (the reward signal used during training) continues to rise.
This finding underscores that even in the DPO framework, which does not use an explicit reward model,
the tension between optimization and generalization persists.

\subsubsection{2025: New Advances in RL Efficiency and Reasoning Distillation}
\label{sec:t7-rl-distill-2025}

\paragraph{RLVR as compute-efficient RL.}
RLVR methods such as DAPO \cite{yu2025dapo} and GSPO \cite{zheng2025gspo}
(\crossref{2}{sec:t2-rlvr}) have not only achieved breakthroughs in reasoning performance
but also provided new answers to the efficiency challenges of RL-based post-training.
DAPO further simplified the GRPO training pipeline by removing the KL constraint and the value model,
making RLVR training lower in both compute and memory overhead compared to traditional PPO-based RLHF.
This trend suggests that: \textbf{for verifiable tasks, RLVR is simultaneously the most effective and most efficient RL method}.

\paragraph{Scaling up reasoning distillation.}
After DeepSeek-R1 \cite{deepseek2025r1} confirmed the distillability of reasoning capabilities,
reasoning distillation rapidly became the most economical pathway for acquiring reasoning abilities.
R1 distilled models (e.g., R1-Distill-Qwen-32B) require only standard SFT training
to significantly outperform non-distilled models of comparable scale on reasoning benchmarks.
Marco-o1 v2 \cite{yin2025marcoo1v2} further explored methods for expanding the distillation bottleneck.
The implication for the efficiency domain is that:
\textbf{``large model RL training $\to$ small model SFT distillation'' may be the optimal compute allocation strategy for acquiring reasoning capabilities}.

\paragraph{Adjustable reasoning budgets.}
The long2short strategy of Kimi k1.5 \cite{moonshot2025kimi} and
the thinking mode fusion of Qwen3 \cite{yang2025qwen3}
have opened the research direction of ``adjustable reasoning budgets'':
models can dynamically adjust inference-time compute based on problem difficulty,
responding quickly to simple problems and thinking deeply on complex ones.
This extends Thread~7's efficiency optimization from training-time to inference-time.

\subsection{Open Problems}
\label{sec:t7-open}

\begin{openproblembox}
\textbf{Model Collapse: The Recursive Trap of Synthetic Data.}\\
Shumailov et al.\ \cite{shumailov2023modelcollapse} revealed a fundamental risk:
when models are repeatedly trained on data generated by themselves or predecessor models,
the output distribution progressively degrades (mode collapse),
and low-frequency but valuable information is irreversibly lost.
As the proportion of synthetic data in post-training continues to rise
(Self-Instruct, UltraChat, and WizardLM's Evol-Instruct all rely on model-generated data),
the risk of model collapse grows increasingly severe.
How to maintain output distribution diversity while using synthetic data at scale---whether through Molmo-style purely human-annotated data as an ``anchor,'' or through theoretically guided synthetic-real data mixing ratios---remains an open problem.
\end{openproblembox}

\begin{openproblembox}
\textbf{Optimal Composition of Synthetic Data.}
The optimal balance among the three dimensions of synthetic data---quality, diversity, and complexity---remains unclear.
Phi-1's ``textbook quality'' strategy, WizardLM's ``complexity evolution'' strategy, and LIMA's ``less is more'' strategy each emphasize different data characteristics---yet no unified theoretical framework currently exists to guide the optimal generation strategy for synthetic data.
\end{openproblembox}

\begin{openproblembox}
\textbf{Scaling Laws for RL-Based Post-Training.}
GRPO and DeepSeek-R1 have demonstrated the enormous potential of RL-based post-training for reasoning,
but its scaling behavior has not yet been systematically characterized.
Key questions include:
What is the optimal allocation strategy for RL training compute?
How should the group size $G$ scale with model size and task complexity?
How does the coverage range of verifiable rewards affect scaling behavior?
\end{openproblembox}

\begin{openproblembox}
\textbf{Unification of Data Quality Metrics.}
Existing data selection methods (IFD, DEITA, LESS, etc.) each define different quality metrics,
and the correlations and complementarities among these metrics have not been thoroughly studied.
A universal, composable data quality metric would greatly advance the systematic development of data engineering.
\end{openproblembox}

\begin{openproblembox}
\textbf{Capability Boundaries of Distillation vs.\ Native Training.}
Zephyr and Orca have demonstrated the powerful potential of distillation,
but does a theoretical upper bound exist on the capabilities of distilled models?
In other words, can a student ever surpass the teacher in capability dimensions where the teacher itself is weak?
Preliminary explorations in weak-to-strong generalization \cite{burns2023weak} suggest that
students can surpass teachers in certain cases,
but systematic theoretical analysis remains lacking.
\crossref{3}{sec:t3}
\end{openproblembox}

\begin{openproblembox}
\textbf{Adaptive Reasoning Budget Optimization.}
Thinking models have introduced a new efficiency dimension: adaptive allocation of inference-time compute.
Current methods (e.g., s1's budget forcing \cite{muennighoff2025s1}, Qwen3's thinking mode fusion)
provide coarse-grained reasoning budget control, but the optimal adaptive strategy remains unclear.
Key questions include: How can models be trained during post-training to autonomously assess problem difficulty?
How can token-level rather than sequence-level fine-grained control of reasoning budgets be achieved?
Can reasoning budget optimization be jointly optimized with PEFT and distillation techniques?
\end{openproblembox}

\begin{openproblembox}
\textbf{A Unified Benchmark for Post-Training Efficiency.}
No standardized benchmark currently exists for fairly comparing the cost-effectiveness of different efficiency techniques.
An ideal benchmark should simultaneously measure FLOPs, peak memory, wall-clock time,
final task performance, and data annotation cost,
enabling researchers to compare the ROI of different strategy combinations---such as LoRA vs.\ full FT,
synthetic data vs.\ human data, and PPO vs.\ GRPO---within a unified framework.
\end{openproblembox}

\subsection{Chapter Summary}

\begin{figure}[!ht]
\centering
\begin{tikzpicture}[
  node distance=0.6cm and 1.2cm,
  every node/.style={font=\small},
  milestone/.style={rectangle, rounded corners, draw=thread7, fill=thread7!10,
    text width=4.2cm, minimum height=0.8cm, align=center, font=\small\bfseries},
  branch/.style={rectangle, rounded corners, draw=gray!60, fill=gray!5,
    text width=3.2cm, minimum height=0.7cm, align=center, font=\small},
  arrow/.style={-{Stealth[length=2.5mm]}, thick, thread7!70},
  brancharrow/.style={-{Stealth[length=2mm]}, gray!60, thick},
]
% Main timeline
\node[milestone] (lora) {LoRA (2022)\\Parameter-Efficient Fine-Tuning};
\node[milestone, below=1.2cm of lora] (lima) {LIMA (2023)\\1000 Samples Suffice};
\node[milestone, below=1.2cm of lima] (qlora) {QLoRA (2023)\\4-bit Quantized Adaptation};

% Branches
\node[branch, below left=1.2cm and 0.5cm of qlora] (brA) {Branch A: Synthetic Data\\Self-Instruct, Alpaca};
\node[branch, below right=1.2cm and 0.5cm of qlora] (brB) {Branch B: Data Curation\\Phi-1, IFD, DEITA};

% Convergence
\node[milestone, below=2.8cm of qlora] (zephyr) {Zephyr (2023)\\Distillation Pipeline};
\node[milestone, below=1.2cm of zephyr] (grpo) {GRPO (2024)\\Algorithmic Simplification};
\node[branch, below=1.2cm of grpo] (conv) {Convergence: Data + Parameter + Algorithm\\Three-Way Efficiency Fusion};

% Arrows
\draw[arrow] (lora) -- (lima);
\draw[arrow] (lima) -- (qlora);
\draw[brancharrow] (qlora) -- (brA);
\draw[brancharrow] (qlora) -- (brB);
\draw[arrow] (qlora) -- (zephyr);
\draw[arrow] (zephyr) -- (grpo);
\draw[brancharrow] (brA) -- (zephyr);
\draw[brancharrow] (brB) -- (zephyr);
\draw[arrow] (grpo) -- (conv);
\end{tikzpicture}
\caption{Thread 7 evolution path: from parameter efficiency breakthroughs to three-way fusion of data, parameter, and algorithm efficiency.
Brown nodes denote landmark works; gray nodes denote important branches.}
\label{fig:efficiency-evolution}
\end{figure}

\begin{table}[!ht]
\centering
\caption{Comparison of core methods in Thread 7. The Efficiency Dimension distinguishes parameter efficiency, data efficiency, and algorithmic efficiency.}
\label{tab:efficiency-comparison}
\small
\begin{tabularx}{\textwidth}{l >{\raggedright\arraybackslash}p{2.2cm} >{\raggedright\arraybackslash}p{2.2cm} >{\raggedright\arraybackslash}X}
\toprule
\textbf{Method} & \textbf{Efficiency Dim.} & \textbf{Resource Savings} & \textbf{Key Innovation} \\
\midrule
LoRA & Parameter & 0.1\% trainable params & Low-rank adapter matrices; frozen backbone \\
QLoRA & Parameter & 4-bit base + LoRA & NF4 quantization; paged optimizers \\
LIMA & Data & 1K samples suffice & Superficial Alignment Hypothesis \\
Self-Instruct & Data & \$600 API cost & Bootstrap 52K instructions from 175 seeds \\
Zephyr & Data + Algorithm & No human annotation & GPT-4 distillation + DPO; 7B $>$ 70B \\
GRPO & Algorithm & No value network & Group-normalized advantage replaces critic \\
\bottomrule
\end{tabularx}
\end{table}

The evolution of Thread~7 reveals the inevitable trend in post-training from ``brute-force scaling'' to ``precision engineering'':

\begin{enumerate}[leftmargin=2em]
  \item \textbf{The data efficiency revolution} (2023):
    LIMA demonstrated that 1,000 carefully selected samples can rival the SFT performance of tens of thousands of data points.
    AlpaGasus and IFD established methodologies for data quality assessment.
    Phi-1's ``textbook-quality data'' approach pushed data curation to its extreme.

  \item \textbf{Parameter efficiency breakthroughs} (2023--2024):
    LoRA compressed trainable parameters to below 0.1\%.
    QLoRA enabled efficient fine-tuning under 4-bit quantization.
    DoRA further separated directional and magnitude adaptation.

  \item \textbf{Algorithmic efficiency optimization} (2024--2025):
    GRPO and REINFORCE++ simplified PPO's four-model architecture to two- or single-model designs.
    ReMax completely eliminated the critic network.
    Distillation techniques enabled small models to acquire the reasoning capabilities of large models.

  \item \textbf{System-level integration} (2025--2026):
    The three efficiency pathways---data, parameter, and algorithm---have begun to converge.
    The combination of efficient algorithms + PEFT + curated data has become the standard practice.
    Model merging techniques provide a new paradigm for combining capabilities without additional training.
\end{enumerate}

The core insight from efficiency research is that the bottleneck of post-training is often not compute volume, but data quality and algorithm design.
This thread is deeply intertwined with Thread~1 (foundational methods of data curation, \crossref{1}{sec:t1}),
Thread~2 (simplification of preference optimization algorithms, \crossref{2}{sec:t2}), and
Thread~4 (distillation of reasoning capabilities, \crossref{4}{sec:t4}).

\paragraph{Preview of the next chapter.}
With this, all seven evolutionary threads have been fully presented.
The final chapter will return to a global perspective, synthesizing the cross-thread convergence trends,
distilling the macro-level evolutionary patterns of the post-training field, and outlining key future research directions.
